% GENERATED FILE. Edit canonical lessons, then run npm run generate:books.
% canonical-source-hash: fnv1a64:8acd75b607940591
% canonical-lessons: TE-C487-veccincu, TE-C487-ada-ceyu, TE-C487-giraki, TE-C487-labhincu, TE-C487-ammudupovu

\chapter{Spending and Saving}
\label{ch:te-spending-and-saving}
\hlchaptermodality{487}

\begin{chapteropening}
\textbf{By the end of this chapter:} \emph{I can say to spend money, to save money, demand, to be available and to sell out.}

Complete the last lesson of chapter 487: I can say to spend money, to save money, demand, to be available and to sell out.
\end{chapteropening}

\section[\texorpdfstring{\te{వెచ్చించు} (vecciñcu)}{వెచ్చించు (vecciñcu)}]{\te{వెచ్చించు} (vecciñcu) — to spend money}

\label{lesson:TE-C487-veccincu}



\begin{quote}
Before the new one: say the Telugu for value, then the Telugu for to weigh.
\end{quote}

\subsection*{You'll want to know: \te{వెచ్చించు}}
\textbf{\te{వెచ్చించు}} — \emph{vecciñcu} — \textquotedblleft{}to spend money\textquotedblright{}.

\begin{quote}
\te{ఇల్లు} \te{కట్టడానికి} \te{చాలా} \te{డబ్బు} \te{వెచ్చించాం}.
\end{quote}

\begin{quote}
\emph{illu kaṭṭaḍāniki cālā ḍabbu vecciñcāṁ.}
\end{quote}

\begin{quote}
We spent a lot of money building the house.
\end{quote}

Past (I …): \te{వెచ్చించాను}, \emph{vecciñcānu}.

A little more formal than \te{ఖర్చు} \te{చేయు}. Time can be spent this way too: \te{సమయం} \te{వెచ్చించు}.

\begin{cousinweb}
\textbf{\te{వెచ్చించు}} (\emph{vecciñcu}) is a Telugu verb in -\te{ఇంచు} whose origin is not settled.
\end{cousinweb}

\begin{grammarlens}[title={What you've built}]
One more word for getting around.
\end{grammarlens}

\subsection*{Your turn}
\begin{itemize}
\raggedright
  \item \emph{Say it:} \emph{vecciñcu}
  \item \emph{Say it:} \emph{vecciñcu}, once more
  \item \emph{Say it:} \emph{illu kaṭṭaḍāniki cālā ḍabbu vecciñcāṁ.} — We spent a lot of money building the house.
  \item \emph{Recall:} say the Telugu for value, then the Telugu for to weigh, then say \emph{vecciñcu} again
\end{itemize}

\begin{tcolorbox}[breakable,colback=teal!4,colframe=teal!35!black,title={Before you move on}]
What does \te{వెచ్చించు} mean? (\textquotedblleft{}To spend money\textquotedblright{}.) Say it once more.
\end{tcolorbox}
\section[\texorpdfstring{\te{ఆదా} \te{చేయు} (ādā cēyu)}{ఆదా చేయు (ādā cēyu)}]{\te{ఆదా} \te{చేయు} (ādā cēyu) — to save money}

\label{lesson:TE-C487-ada-ceyu}



\begin{quote}
Before the new one: say the Telugu for to weigh, then the Telugu for to spend money.
\end{quote}

\subsection*{You'll want to know: \te{ఆదా} \te{చేయు}}
\textbf{\te{ఆదా} \te{చేయు}} — \emph{ādā cēyu} — \textquotedblleft{}to save money\textquotedblright{}.

\begin{quote}
\te{బస్సులో} \te{వెళ్ళి} \te{కిరాయి} \te{ఆదా} \te{చేశాం}.
\end{quote}

\begin{quote}
\emph{bassulō veḷḷi kirāyi ādā cēśāṁ.}
\end{quote}

\begin{quote}
We went by bus and saved on the fare.
\end{quote}

Past (I …): \te{ఆదా} \te{చేశాను}, \emph{ādā cēśānu}.

Said of money, time or fuel saved: \te{సమయం} \te{ఆదా} \te{చేయు}.

\begin{cousinweb}
\textbf{\te{ఆదా} \te{చేయు}} (\emph{ādā cēyu}) is a Telugu phrase: \te{ఆదా}, \textquotedblleft{}saving\textquotedblright{}, whose origin is not settled, and the native \te{చేయు}, \textquotedblleft{}to do\textquotedblright{} (DEDR 1957).
\end{cousinweb}

\begin{grammarlens}[title={What you've built}]
One more word for getting around.
\end{grammarlens}

\subsection*{Your turn}
\begin{itemize}
\raggedright
  \item \emph{Say it:} \emph{ādā cēyu}
  \item \emph{Say it:} \emph{ādā cēyu}, once more
  \item \emph{Say it:} \emph{bassulō veḷḷi kirāyi ādā cēśāṁ.} — We went by bus and saved on the fare.
  \item \emph{Recall:} say the Telugu for to weigh, then the Telugu for to spend money, then say \emph{ādā cēyu} again
\end{itemize}

\begin{tcolorbox}[breakable,colback=teal!4,colframe=teal!35!black,title={Before you move on}]
What does \te{ఆదా} \te{చేయు} mean? (\textquotedblleft{}To save money\textquotedblright{}.) Say it once more.
\end{tcolorbox}
\section[\texorpdfstring{\te{గిరాకీ} (girākī)}{గిరాకీ (girākī)}]{\te{గిరాకీ} (girākī) — demand}

\label{lesson:TE-C487-giraki}



\begin{quote}
Before the new one: say the Telugu for to spend money, then the Telugu for to save money.
\end{quote}

\subsection*{You'll want to know: \te{గిరాకీ}}
\textbf{\te{గిరాకీ}} — \emph{girākī} — \textquotedblleft{}demand\textquotedblright{}.

\begin{quote}
\te{వేసవిలో} \te{మామిడి} \te{పండ్లకు} \te{గిరాకీ} \te{ఎక్కువ}.
\end{quote}

\begin{quote}
\emph{vēsavilō māmiḍi paṇḍlaku girākī ekkuva.}
\end{quote}

\begin{quote}
In summer there is high demand for mangoes.
\end{quote}

Demand for goods in a market; also customers for a shop: \te{గిరాకీ} \te{లేదు}, there are no buyers.

\begin{cousinweb}
\textbf{\te{గిరాకీ}} (\emph{girākī}) is a market word for demand whose origin is not settled.
\end{cousinweb}

\begin{grammarlens}[title={What you've built}]
One more word for getting around.
\end{grammarlens}

\subsection*{Your turn}
\begin{itemize}
\raggedright
  \item \emph{Say it:} \emph{girākī}
  \item \emph{Say it:} \emph{girākī}, once more
  \item \emph{Say it:} \emph{vēsavilō māmiḍi paṇḍlaku girākī ekkuva.} — In summer there is high demand for mangoes.
  \item \emph{Recall:} say the Telugu for to spend money, then the Telugu for to save money, then say \emph{girākī} again
\end{itemize}

\begin{tcolorbox}[breakable,colback=teal!4,colframe=teal!35!black,title={Before you move on}]
What does \te{గిరాకీ} mean? (\textquotedblleft{}Demand\textquotedblright{}.) Say it once more.
\end{tcolorbox}
\section[\texorpdfstring{\te{లభించు} (labhiñcu)}{లభించు (labhiñcu)}]{\te{లభించు} (labhiñcu) — to be available; to be obtained}

\label{lesson:TE-C487-labhincu}



\begin{quote}
Before the new one: say the Telugu for to save money, then the Telugu for demand.
\end{quote}

\subsection*{You'll want to know: \te{లభించు}}
\textbf{\te{లభించు}} — \emph{labhiñcu} — \textquotedblleft{}to be available; to be obtained\textquotedblright{}.

\begin{quote}
\te{ఆ} \te{పుస్తకం} \te{మా} \te{పట్టణంలో} \te{లభించదు}.
\end{quote}

\begin{quote}
\emph{ā pustakaṁ mā paṭṭaṇaṁlō labhiñcadu.}
\end{quote}

\begin{quote}
That book is not available in our town.
\end{quote}

Past (it …): \te{లభించింది}, \emph{labhiñcindi}.

The person who gets it takes -\te{కి}: \te{నాకు} \te{ఉద్యోగం} \te{లభించింది}, I got a job. Its subject is a thing, so the past shown is the third person: \te{లభించింది}, it was got.

\begin{cousinweb}
\textbf{\te{లభించు}} (\emph{labhiñcu}) is built from Sanskrit \emph{labh}, \textquotedblleft{}to take, to obtain\textquotedblright{} (Monier-Williams), with the Telugu verb ending -\te{ఇంచు}; \te{లాభం}, a profit, is the same root.
\end{cousinweb}

\begin{grammarlens}[title={What you've built}]
One more word for getting around.
\end{grammarlens}

\subsection*{Your turn}
\begin{itemize}
\raggedright
  \item \emph{Say it:} \emph{labhiñcu}
  \item \emph{Say it:} \emph{labhiñcu}, once more
  \item \emph{Say it:} \emph{ā pustakaṁ mā paṭṭaṇaṁlō labhiñcadu.} — That book is not available in our town.
  \item \emph{Recall:} say the Telugu for to save money, then the Telugu for demand, then say \emph{labhiñcu} again
\end{itemize}

\begin{tcolorbox}[breakable,colback=teal!4,colframe=teal!35!black,title={Before you move on}]
What does \te{లభించు} mean? (\textquotedblleft{}To be available; to be obtained\textquotedblright{}.) Say it once more.
\end{tcolorbox}
\section[\texorpdfstring{\te{అమ్ముడుపోవు} (ammuḍupōvu)}{అమ్ముడుపోవు (ammuḍupōvu)}]{\te{అమ్ముడుపోవు} (ammuḍupōvu) — to sell out; to be sold}

\label{lesson:TE-C487-ammudupovu}



\begin{quote}
Before the new one: say the Telugu for demand, then the Telugu for to be available.
\end{quote}

\subsection*{You'll want to know: \te{అమ్ముడుపోవు}}
\textbf{\te{అమ్ముడుపోవు}} — \emph{ammuḍupōvu} — \textquotedblleft{}to sell out; to be sold\textquotedblright{}.

\begin{quote}
\te{పండుగ} \te{ముందు} \te{టిక్కెట్లన్నీ} \te{అమ్ముడుపోయాయి}.
\end{quote}

\begin{quote}
\emph{paṇḍuga mundu ṭikkeṭlannī ammuḍupōyāyi.}
\end{quote}

\begin{quote}
Before the festival all the tickets sold out.
\end{quote}

Past (it …): \te{అమ్ముడుపోయింది}, \emph{ammuḍupōyindi}.

\te{అమ్ముడు}, selling, + \te{పోవు}. Said of goods; said of a person, it means to be bought off. Its subject is a thing, so the past shown is the third person: \te{అమ్ముడుపోయింది}, it sold out.

\begin{cousinweb}
\textbf{\te{అమ్ముడుపోవు}} (\emph{ammuḍupōvu}) is native Dravidian: DEDR 186 lists Telugu \emph{ammuḍu}, \textquotedblleft{}sale, selling\textquotedblright{}, beside \emph{ammu}, \textquotedblleft{}to sell\textquotedblright{}, here with \te{పోవు}, \textquotedblleft{}to go\textquotedblright{} (DEDR 4572).
\end{cousinweb}

\begin{grammarlens}[title={What you've built}]
One more word for getting around.
\end{grammarlens}

\subsection*{Your turn}
\begin{itemize}
\raggedright
  \item \emph{Say it:} \emph{ammuḍupōvu}
  \item \emph{Say it:} \emph{ammuḍupōvu}, once more
  \item \emph{Say it:} \emph{paṇḍuga mundu ṭikkeṭlannī ammuḍupōyāyi.} — Before the festival all the tickets sold out.
  \item \emph{Recall:} say the Telugu for demand, then the Telugu for to be available, then say \emph{ammuḍupōvu} again
\end{itemize}

\begin{tcolorbox}[breakable,colback=teal!4,colframe=teal!35!black,title={Before you move on}]
What does \te{అమ్ముడుపోవు} mean? (\textquotedblleft{}To sell out; to be sold\textquotedblright{}.) Say it once more.
\end{tcolorbox}
